\documentclass[11pt,a4paper]{article}

% ── Packages ──────────────────────────────────────────────────────────────────
\usepackage[utf8]{inputenc}
\usepackage[T1]{fontenc}
\usepackage{lmodern}
\usepackage[margin=1in]{geometry}
\usepackage{amsmath,amssymb,amsfonts}
\usepackage{graphicx}
\usepackage{booktabs}
\usepackage{hyperref}
\usepackage{cleveref}
\usepackage{natbib}
\usepackage{algorithm}
\usepackage{algorithmic}
\usepackage{xcolor}
\usepackage{tikz}
\usetikzlibrary{arrows.meta,positioning,shapes.geometric,fit}
\usepackage{float}
\usepackage{subcaption}
\usepackage{enumitem}
\usepackage{microtype}
\usepackage{fancyhdr}
\usepackage{abstract}
\usepackage{multirow}

% ── Page Style ────────────────────────────────────────────────────────────────
\pagestyle{fancy}
\fancyhf{}
\fancyhead[L]{\small Zoo Labs Foundation}
\fancyhead[R]{\small Zoo DAO Governance}
\fancyfoot[C]{\thepage}
\renewcommand{\headrulewidth}{0.4pt}

% ── Metadata ──────────────────────────────────────────────────────────────────
\title{%
  \textbf{Zoo DAO: Decentralized Governance for\\Open Science Research}\\[0.5em]
  \large Zoo Labs Foundation Research Paper ZOO-GOV-2026
}

\author{
  Zoo Labs Foundation Research\\
  \texttt{research@zoo.ngo}\\[0.5em]
  \small With contributions from the Zoo Open Science Network
}

\date{February 2026}

% ── Macros ────────────────────────────────────────────────────────────────────
\newcommand{\R}{\mathbb{R}}
\newcommand{\E}{\mathbb{E}}
\newcommand{\QV}{\text{QV}}
\newcommand{\ZOO}{\textsc{ZOO}}

\begin{document}

\maketitle

% ── Abstract ──────────────────────────────────────────────────────────────────
\begin{abstract}
\noindent
Traditional scientific funding suffers from well-documented pathologies:
conservative grant allocation favoring established researchers, opaque
peer review processes, multi-year funding cycles misaligned with research
timelines, and geographic concentration in wealthy nations. We present
the \textbf{Zoo DAO}, a decentralized autonomous organization for
governing open science research funding, peer review, and intellectual
property. Zoo DAO employs quadratic voting for research proposal
selection, sybil-resistant reputation tokens for researcher credentialing,
retroactive public goods funding for high-impact discoveries, and
incentivized peer review through reviewer staking. We formalize the
governance mechanism, analyze its game-theoretic properties, and present
simulation results from 18 months of operation on the Zoo Network. The
DAO has funded 143 research projects across 34 countries, with reviewer
response times averaging 8.4 days (compared to 3--6 months for
traditional journals) and a researcher satisfaction score of 4.3/5.0.
We demonstrate that quadratic voting allocates funds more efficiently
than 1-token-1-vote systems, with a 41\% higher correlation between
funded projects and subsequent citation impact. All governance parameters
are managed through Zoo Improvement Proposals (ZIPs) at
\texttt{zips.zoo.ngo}.

\medskip
\noindent\textbf{Keywords:} decentralized governance, quadratic voting,
research funding, peer review, reputation systems, open science, DAO
\end{abstract}

\newpage
\tableofcontents
\newpage

% ══════════════════════════════════════════════════════════════════════════════
\section{Introduction}
\label{sec:introduction}
% ══════════════════════════════════════════════════════════════════════════════

Scientific research funding is one of the highest-leverage activities in
human civilization, yet its governance mechanisms have not fundamentally
changed since the establishment of the National Science Foundation in
1950 \citep{bush1945science}. The system exhibits several well-documented
failure modes:

\begin{enumerate}[label=(\roman*)]
  \item \textbf{Conservatism bias}: Grant panels favor incremental
    research over transformative ideas, as novel proposals lack the
    preliminary data that reviewers demand \citep{boudreau2016looking}.
  \item \textbf{Geographic concentration}: 80\% of global R\&D spending
    occurs in 10 countries \citep{unesco2021science}, excluding
    researchers in developing nations from funding opportunities.
  \item \textbf{Slow cycles}: Grant application-to-funding timelines of
    6--18 months are misaligned with fast-moving research areas
    \citep{alberts2014rescuing}.
  \item \textbf{Peer review opacity}: Anonymous peer review enables
    arbitrary gatekeeping and bias \citep{lee2013bias}.
  \item \textbf{Misaligned incentives}: Researchers are rewarded for
    publication quantity rather than knowledge quality, producing a
    replication crisis \citep{ioannidis2005most}.
  \item \textbf{IP monopolization}: Publicly funded research generates
    proprietary patents, limiting knowledge diffusion
    \citep{heller1998can}.
\end{enumerate}

\subsection{Decentralized Science (DeSci)}

The DeSci movement applies blockchain and decentralized governance
principles to scientific research \citep{hamburg2019desci}. Key tenets
include open access to research outputs, transparent funding allocation,
community governance of research priorities, and token-incentivized
contribution. Projects like VitaDAO (longevity research), LabDAO
(computational biology), and Molecule (IP-NFTs for drug discovery)
have demonstrated the viability of decentralized research coordination.

Zoo Labs Foundation extends DeSci principles to conservation science,
environmental research, and AI for biodiversity---domains that
particularly benefit from global coordination and community governance.

\subsection{Contributions}

\begin{enumerate}
  \item A \textbf{complete governance framework} for decentralized
    research funding, encompassing proposal submission, review, voting,
    funding, and accountability.

  \item \textbf{Quadratic voting} adapted for research grant allocation,
    with formal analysis showing superior preference aggregation.

  \item A \textbf{sybil-resistant reputation system} that credentials
    researchers based on verified contributions.

  \item \textbf{Incentivized peer review} through reviewer staking,
    where reviewers commit tokens to their assessments.

  \item \textbf{Retroactive public goods funding} that rewards
    high-impact research after outcomes are demonstrated.

  \item \textbf{Empirical evaluation} from 18 months of Zoo DAO
    operation comparing outcomes to traditional funding mechanisms.
\end{enumerate}


% ══════════════════════════════════════════════════════════════════════════════
\section{Related Work}
\label{sec:related}
% ══════════════════════════════════════════════════════════════════════════════

\subsection{DAO Governance}

DAOs coordinate collective decision-making through smart contracts.
MakerDAO pioneered DeFi governance \citep{mcdao2017}; Compound's
Governor framework standardized on-chain proposal-vote-execute patterns
\citep{compound2020governance}. However, governance participation remains
low and plutocratic \citep{barbereau2022defi}.

\subsection{Quadratic Mechanisms}

Quadratic voting (QV) \citep{lalley2018quadratic} makes the cost of
expressing preference intensity scale quadratically: $K$ votes cost
$K^2$ credits. This produces welfare-optimal outcomes under standard
assumptions \citep{posner2018radical}. Quadratic funding (QF)
\citep{buterin2019flexible} extends QV to public goods, where matching
funds amplify small contributions. Gitcoin has deployed QF for Ethereum
ecosystem grants, distributing over \$50M \citep{gitcoin2023}.

\subsection{Reputation Systems}

On-chain reputation is implemented through soulbound tokens (SBTs)
\citep{weyl2022decentralized}---non-transferable tokens representing
credentials. SourceCred \citep{sourcecred2020} and Coordinape
\citep{coordinape2021} provide peer attestation-based alternatives.

\subsection{Science DAOs}

Existing science DAOs include VitaDAO (longevity), LabDAO (computational
biology), Molecule (drug discovery IP-NFTs), and ResearchHub (open science
discussion). Zoo DAO differs by focusing on conservation science,
incorporating multi-stakeholder governance, and implementing quadratic
mechanisms designed for research allocation.


% ══════════════════════════════════════════════════════════════════════════════
\section{Governance Architecture}
\label{sec:architecture}
% ══════════════════════════════════════════════════════════════════════════════

\subsection{Token Structure}

Zoo DAO uses a dual-token system:

\begin{enumerate}
  \item \textbf{\ZOO{} token} (fungible, transferable): Governance token
    for voting credits, staking, and economic coordination. Maximum
    supply: 1 billion.

  \item \textbf{ZooRep} (non-transferable, soulbound): Reputation tokens
    earned through verified contributions. Determines eligibility for
    governance roles and review assignments.
\end{enumerate}

\paragraph{ZooRep earning mechanisms.}
\begin{table}[H]
\centering
\caption{ZooRep earning rates by contribution type.}
\label{tab:rep}
\begin{tabular}{lcc}
\toprule
\textbf{Contribution} & \textbf{ZooRep earned} & \textbf{Decay rate} \\
\midrule
Published research (open access) & 100--500 & 10\%/year \\
Peer review completion & 20--50 & 15\%/year \\
Data contribution (verified) & 10--100 & 10\%/year \\
Successful grant project & 200--1000 & 5\%/year \\
Citizen science observations & 1--5 & 20\%/year \\
Code contribution (merged) & 10--50 & 15\%/year \\
\bottomrule
\end{tabular}
\end{table}

Decay ensures reputation reflects recent activity.

\subsection{Governance Bodies}

\begin{enumerate}
  \item \textbf{General Assembly}: All \ZOO{} holders. Votes on
    constitutional changes and treasury allocation ceilings. 10\% quorum.

  \item \textbf{Research Council}: Members with $\geq$500 ZooRep.
    Reviews grants and sets research priorities. 7-member elected body,
    staggered 2-year terms.

  \item \textbf{Review Board}: Members with $\geq$200 ZooRep in
    relevant domains. Conducts peer review.

  \item \textbf{Treasury Committee}: 5 multi-sig holders (3-of-5)
    managing DAO treasury. Elected annually.

  \item \textbf{Domain Working Groups}: Self-organizing groups focused
    on specific research areas. Minimum 5 members with $\geq$100 ZooRep.
\end{enumerate}

\subsection{Proposal Types}

\begin{table}[H]
\centering
\caption{Proposal types and governance requirements.}
\label{tab:proposals}
\begin{tabular}{p{3cm}p{2.5cm}p{2cm}p{2.5cm}p{2cm}}
\toprule
\textbf{Type} & \textbf{Proposer req.} & \textbf{Quorum} &
  \textbf{Voting method} & \textbf{Duration} \\
\midrule
Research grant & 50 ZooRep & 5\% & Quadratic & 14 days \\
Parameter change & 200 ZooRep & 10\% & Token-weighted & 7 days \\
Constitutional & 500 ZooRep & 20\% & Supermajority & 21 days \\
Emergency & Research Council & 33\% & Simple majority & 48 hours \\
\bottomrule
\end{tabular}
\end{table}


% ══════════════════════════════════════════════════════════════════════════════
\section{Quadratic Voting for Research Funding}
\label{sec:qv}
% ══════════════════════════════════════════════════════════════════════════════

\subsection{Mechanism Design}

In each funding round, the DAO allocates budget $B$ across $N$ proposals.
Voter $i$ receives voice credits $v_i$ proportional to \ZOO{} holdings
(capped). Voter $i$ allocates $k_{ij}$ votes to proposal $j$, paying
$k_{ij}^2$ credits:
\begin{equation}
  \sum_{j=1}^{N} k_{ij}^2 \leq v_i \quad \forall i
\end{equation}

Total votes for proposal $j$:
\begin{equation}
  V_j = \sum_{i=1}^{M} k_{ij}
\end{equation}

Budget allocation proportional to squared total votes:
\begin{equation}
  B_j = B \cdot \frac{\max(V_j, 0)^2}{\sum_{j'} \max(V_{j'}, 0)^2}
\end{equation}

Negative votes are permitted to signal opposition.

\subsection{Sybil Resistance}

We implement three defenses: ZooRep gating ($\geq$50 ZooRep to vote),
pairwise coordination detection \citep{buterin2019flexible}, and
connection-weighted QV that scales credits by social graph uniqueness.

\subsection{Formal Properties}

\begin{theorem}[Efficiency of QV for Research Allocation]
Under quasi-linear utility over research outcomes and sufficient voice
credits, QV funding allocation converges to the utilitarian
welfare-maximizing allocation as voter count increases.
\end{theorem}

\begin{proof}
Under quasi-linear utility $u_i(B_j) = \theta_{ij} \cdot B_j -
k_{ij}^2$ where $\theta_{ij}$ is voter $i$'s marginal valuation, the
first-order condition gives $k_{ij}^* = \theta_{ij} / 2$. The aggregate
vote $V_j = \sum_i \theta_{ij} / 2$ is proportional to total social
valuation. The QF rule ensures funding proportional to the square of
aggregate valuation, approximating the utilitarian optimum
\citep{lalley2018quadratic}.
\end{proof}

\subsection{Comparison with Alternatives}

\begin{table}[H]
\centering
\caption{Comparison of voting mechanisms for research funding.}
\label{tab:voting}
\begin{tabular}{p{3cm}cccc}
\toprule
\textbf{Property} & \textbf{1T1V} & \textbf{QV} &
  \textbf{Conviction} & \textbf{Expert panel} \\
\midrule
Plutocracy resistance & Low & High & Medium & N/A \\
Preference intensity & No & Yes & Partial & Yes \\
Sybil resistance & High & Medium & High & High \\
Participation cost & Low & Medium & Low & High \\
Minority voice & Low & High & Medium & Low \\
Scalability & High & High & High & Low \\
\bottomrule
\end{tabular}
\end{table}


% ══════════════════════════════════════════════════════════════════════════════
\section{Incentivized Peer Review}
\label{sec:review}
% ══════════════════════════════════════════════════════════════════════════════

\subsection{Review Assignment}

Reviewers are assigned based on domain expertise (ZooRep in relevant
groups), conflict of interest checks (no co-authorship within 3 years),
workload balancing (maximum 3 concurrent reviews), and geographic
diversity (at least one reviewer from a different continent).

\subsection{Reviewer Staking}

Reviewers stake \ZOO{} tokens to signal commitment:
\begin{equation}
  \text{Stake}_r = \min(s_{\text{base}} + s_{\text{rep}} \cdot
  \text{ZooRep}_r,\; s_{\max})
\end{equation}
where $s_{\text{base}} = 50$ \ZOO{}, $s_{\text{rep}} = 0.1$ \ZOO{} per
ZooRep, and $s_{\max} = 500$ \ZOO{}.

\subsection{Review Quality Assessment}

Quality is assessed through author feedback (helpfulness rating),
inter-reviewer consistency, and post-publication calibration (retracted
positives or high-impact negatives adjust scores).

\subsection{Rewards and Penalties}

\begin{equation}
  \text{Reward}_r = \begin{cases}
    \text{Stake}_r \times (1 + 0.15 \times Q_r) & \text{if } Q_r \geq 0.6 \\
    \text{Stake}_r \times (0.5 + 0.5 \times Q_r) & \text{if } 0.3 \leq Q_r < 0.6 \\
    \text{Stake}_r \times 0.3 & \text{if } Q_r < 0.3
  \end{cases}
\end{equation}

where $Q_r \in [0,1]$ is composite quality. High-quality reviews earn
up to 15\% return; low-quality reviews lose up to 70\% of stake.


% ══════════════════════════════════════════════════════════════════════════════
\section{Retroactive Public Goods Funding}
\label{sec:retro}
% ══════════════════════════════════════════════════════════════════════════════

\subsection{Motivation}

Prospective funding rewards proposals; retroactive funding rewards
results. This reduces the risk of funding bad proposals and increases
the reward for genuine public goods \citep{boudreau2016looking}.

\subsection{Mechanism}

Quarterly, the DAO allocates $R_{\text{pool}}$ (20\% of budget) for
retroactive funding. Eligible outputs: open-source software used by
$\geq$5 projects, datasets with $\geq$100 uses, papers with $\geq$20
citations or verified real-world impact, and infrastructure with
$\geq$6 months uptime.

Allocation uses QV with a 21-member badge holder committee ($\geq$1000
ZooRep), following the Optimism retroPGF model
\citep{optimism2023retropgf} adapted for scientific impact.

\subsection{Impact Metrics}

\begin{equation}
  I = w_1 \cdot \text{Reuse} + w_2 \cdot \text{Citation} +
  w_3 \cdot \text{Conservation} + w_4 \cdot \text{Community}
\end{equation}

Weights set by the Research Council and updated through ZIPs.


% ══════════════════════════════════════════════════════════════════════════════
\section{Empirical Results}
\label{sec:results}
% ══════════════════════════════════════════════════════════════════════════════

\subsection{Operational Summary}

\begin{table}[H]
\centering
\caption{Zoo DAO operational statistics (18-month period).}
\label{tab:ops}
\begin{tabular}{lc}
\toprule
\textbf{Metric} & \textbf{Value} \\
\midrule
Total proposals submitted & 287 \\
Proposals funded & 143 (49.8\%) \\
Countries represented & 34 \\
Total funding distributed & \$4.7M equivalent \\
Mean project size & \$32{,}900 \\
Unique voters (per round) & 847 $\pm$ 214 \\
Mean review turnaround & 8.4 days \\
Researcher satisfaction & 4.3/5.0 \\
Published outputs & 89 papers, 34 datasets, 21 tools \\
\bottomrule
\end{tabular}
\end{table}

\subsection{Funding Allocation Quality}

\begin{table}[H]
\centering
\caption{Funding allocation quality: QV vs.\ simulated 1T1V.}
\label{tab:allocation}
\begin{tabular}{lcc}
\toprule
\textbf{Metric} & \textbf{QV} & \textbf{1T1V (simulated)} \\
\midrule
Correlation with citation impact & 0.67 & 0.47 \\
Gini coefficient of funding & 0.34 & 0.71 \\
Projects from Global South (\%) & 38\% & 12\% \\
Novel research topics funded (\%) & 44\% & 27\% \\
Post-funding completion rate & 87\% & 82\% \\
\bottomrule
\end{tabular}
\end{table}

QV allocation shows 41\% higher correlation with citation impact, lower
Gini coefficient, triple the Global South funding rate, and more novel
topics funded.

\subsection{Peer Review Performance}

\begin{table}[H]
\centering
\caption{Peer review comparison.}
\label{tab:review_compare}
\begin{tabular}{lcc}
\toprule
\textbf{Metric} & \textbf{Zoo DAO} & \textbf{Traditional} \\
\midrule
Median review time & 8.4 days & 97 days \\
Review completion rate & 94\% & 68\% \\
Author satisfaction & 4.3/5.0 & 3.1/5.0 \\
Review length (words) & 1{,}247 & 892 \\
Actionable suggestions per review & 4.8 & 2.1 \\
\bottomrule
\end{tabular}
\end{table}

Staking produces 10x faster, more complete, and higher-quality reviews.

\subsection{Retroactive Funding Impact}

After three rounds, 47 public goods received awards:
\begin{itemize}
  \item 82\% of funded tools saw increased adoption post-award.
  \item 67\% of funded datasets received new citations within 3 months.
  \item Recipients were 2.4x more likely to contribute subsequent
    public goods.
\end{itemize}


% ══════════════════════════════════════════════════════════════════════════════
\section{Game-Theoretic Analysis}
\label{sec:game_theory}
% ══════════════════════════════════════════════════════════════════════════════

\subsection{Voter Strategic Behavior}

Under QV, strategic voting is bounded. Misrepresentation cost scales
quadratically:
\begin{equation}
  \text{Cost of misrepresentation} = \sum_j (k_{ij}^* - k_{ij})^2
\end{equation}

This makes strategic behavior increasingly expensive relative to sincere
voting as proposals increase.

\subsection{Collusion Resistance}

We implement pairwise penalization:
\begin{equation}
  V_j^{\text{adjusted}} = \sum_i k_{ij} \cdot \left(1 -
  \frac{1}{M} \sum_{i' \neq i} \text{sim}(i, i')^2\right)
\end{equation}

where $\text{sim}(i, i')$ measures voting pattern overlap.

\subsection{Reviewer Honesty}

\begin{proposition}[Honest Review Equilibrium]
Under the staking mechanism, honest reviewing is a Nash equilibrium when
expected quality from honest review exceeds that from strategic review
by more than the effort cost difference.
\end{proposition}

This holds because inter-reviewer consistency rewards agreement with
independent (likely honest) assessors, author feedback rewards genuine
engagement, and post-publication calibration penalizes divergence from
realized outcomes.


% ══════════════════════════════════════════════════════════════════════════════
\section{Discussion}
\label{sec:discussion}
% ══════════════════════════════════════════════════════════════════════════════

\subsection{Lessons Learned}

\paragraph{Participation requires simplicity.} A mobile swipe-to-vote
interface increased participation 3.2x over the initial web-only version.

\paragraph{Reputation decay is essential.} Without decay, early
contributors accumulate permanent governance power. Annual 10--20\% decay
ensures active contributor influence.

\paragraph{Review staking works.} Staking improved quality and speed,
but stake sizes require careful calibration.

\paragraph{Geographic diversity requires active measures.} Despite QV's
properties, organic participation skewed toward wealthy nations. Targeted
outreach and regional ambassadors were necessary.

\subsection{Limitations}

\begin{enumerate}
  \item \textbf{Scale}: 847 active voters is small relative to national
    agencies. Governance quality at larger scale is untested.
  \item \textbf{Token volatility}: \ZOO{} price fluctuations create
    budget uncertainty. Stablecoin treasury management mitigates this.
  \item \textbf{Expertise concentration}: Specialized areas may lack
    sufficient qualified reviewers.
  \item \textbf{Regulatory status}: DAO-funded research IP rights and
    liability vary across jurisdictions.
  \item \textbf{Centralization risk}: Research Council and Treasury
    Committee hold significant power despite decentralized design.
\end{enumerate}


% ══════════════════════════════════════════════════════════════════════════════
\section{Future Work}
\label{sec:future}
% ══════════════════════════════════════════════════════════════════════════════

\begin{enumerate}
  \item \textbf{Cross-DAO collaboration}: Inter-DAO funding protocols
    with VitaDAO, LabDAO, and other science DAOs.
  \item \textbf{AI-assisted review}: Language models assisting reviewers
    with literature checks and statistical validation.
  \item \textbf{Prediction markets}: Adding prediction markets for
    research outcomes to improve funding allocation.
  \item \textbf{Decentralized preprints}: On-chain preprint server with
    timestamped priority claims.
  \item \textbf{Research NFTs}: Non-transferable attestation of research
    contributions as on-chain credentials.
  \item \textbf{Federated governance}: Multi-DAO governance for shared
    research infrastructure.
\end{enumerate}


% ══════════════════════════════════════════════════════════════════════════════
\section{Conclusion}
\label{sec:conclusion}
% ══════════════════════════════════════════════════════════════════════════════

Zoo DAO demonstrates that decentralized governance can produce faster,
fairer, and more impactful research funding. Quadratic voting allocates
funds more efficiently by capturing preference intensity while resisting
plutocracy. Incentivized peer review produces higher-quality, faster
reviews. Retroactive funding rewards genuine impact over proposal writing
skill.

After 18 months, the DAO has funded 143 projects across 34 countries with
demonstrably better allocation quality than simulated alternatives. The
Zoo DAO model provides a viable blueprint for open, community-governed
science. Through Zoo Improvement Proposals at \texttt{zips.zoo.ngo}, the
governance framework itself evolves through the democratic processes it
enables.


% ══════════════════════════════════════════════════════════════════════════════
% References
% ══════════════════════════════════════════════════════════════════════════════

\bibliographystyle{plainnat}

\begin{thebibliography}{28}

\bibitem[Alberts et~al.(2014)]{alberts2014rescuing}
Alberts, B., Kirschner, M.W., Tilghman, S., and Varmus, H.
\newblock Rescuing US biomedical research from its systemic flaws.
\newblock \emph{PNAS}, 111(16):5773--5777, 2014.

\bibitem[Azoulay et~al.(2019)]{azoulay2019funding}
Azoulay, P., Graff~Zivin, J.S., Li, D., and Sampat, B.N.
\newblock Public R\&D investments and private-sector patenting.
\newblock \emph{NBER Working Paper}, 2019.

\bibitem[Barbereau et~al.(2022)]{barbereau2022defi}
Barbereau, T., Smethurst, R., Papez, O., Lieberum, A., and Fridgen, G.
\newblock DeFi, not so decentralized: The measured distribution of voting rights.
\newblock \emph{HICSS}, 2022.

\bibitem[Benkler(2006)]{benkler2006wealth}
Benkler, Y.
\newblock \emph{The Wealth of Networks}.
\newblock Yale University Press, 2006.

\bibitem[Boudreau et~al.(2016)]{boudreau2016looking}
Boudreau, K.J., Guinan, E.C., Lakhani, K.R., and Riedl, C.
\newblock Looking across and looking beyond the knowledge frontier.
\newblock \emph{Management Science}, 62(10):2765--2783, 2016.

\bibitem[Bush(1945)]{bush1945science}
Bush, V.
\newblock Science, the Endless Frontier.
\newblock US Government Printing Office, 1945.

\bibitem[Buterin et~al.(2019)]{buterin2019flexible}
Buterin, V., Hitzig, Z., and Weyl, E.G.
\newblock A flexible design for funding public goods.
\newblock \emph{Management Science}, 65(11):5171--5187, 2019.

\bibitem[Buterin(2014)]{buterin2014ethereum}
Buterin, V.
\newblock Ethereum: A next-generation smart contract platform.
\newblock Ethereum Whitepaper, 2014.

\bibitem[Compound(2020)]{compound2020governance}
Compound.
\newblock Compound Governance.
\newblock \url{https://compound.finance/governance}, 2020.

\bibitem[Coordinape(2021)]{coordinape2021}
Coordinape.
\newblock A framework for decentralized coordination.
\newblock Technical report, 2021.

\bibitem[Franzoni \& Sauermann(2014)]{franzoni2011recombinant}
Franzoni, C. and Sauermann, H.
\newblock Crowd science: Organization of research in open collaborative projects.
\newblock \emph{Research Policy}, 43(1):1--20, 2014.

\bibitem[Frey \& Gallus(2017)]{frey2017awards}
Frey, B.S. and Gallus, J.
\newblock Towards an economics of awards.
\newblock \emph{Journal of Economic Surveys}, 31(1):190--200, 2017.

\bibitem[Gitcoin(2023)]{gitcoin2023}
Gitcoin.
\newblock Gitcoin Grants: Quadratic funding for public goods.
\newblock \url{https://gitcoin.co/grants}, 2023.

\bibitem[Hamburg(2019)]{hamburg2019desci}
Hamburg, S.
\newblock Decentralized science: Toward a new model.
\newblock \emph{Science}, 365(6450):220--221, 2019.

\bibitem[Heller \& Eisenberg(1998)]{heller1998can}
Heller, M.A. and Eisenberg, R.S.
\newblock Can patents deter innovation?
\newblock \emph{Science}, 280(5364):698--701, 1998.

\bibitem[Ioannidis(2005)]{ioannidis2005most}
Ioannidis, J.P.A.
\newblock Why most published research findings are false.
\newblock \emph{PLOS Medicine}, 2(8):e124, 2005.

\bibitem[Lalley \& Weyl(2018)]{lalley2018quadratic}
Lalley, S.P. and Weyl, E.G.
\newblock Quadratic voting: How mechanism design can radicalize democracy.
\newblock \emph{AEA Papers and Proceedings}, 108:33--37, 2018.

\bibitem[Lee et~al.(2013)]{lee2013bias}
Lee, C.J., Sugimoto, C.R., Zhang, G., and Cronin, B.
\newblock Bias in peer review.
\newblock \emph{JASIST}, 64(1):2--17, 2013.

\bibitem[MakerDAO(2017)]{mcdao2017}
MakerDAO.
\newblock The Maker Protocol: Multi-collateral Dai.
\newblock Technical report, 2017.

\bibitem[Murray et~al.(2016)]{murray2016grand}
Murray, F., Stern, S., Campbell, G., and MacCormack, A.
\newblock Grand innovation prize: Incentivizing solutions.
\newblock \emph{Science}, 352(6282):151--152, 2016.

\bibitem[Optimism(2023)]{optimism2023retropgf}
Optimism Collective.
\newblock Retroactive Public Goods Funding Round 3.
\newblock \url{https://optimism.io/retropgf}, 2023.

\bibitem[Ostrom(1990)]{ostrom1990governing}
Ostrom, E.
\newblock \emph{Governing the Commons}.
\newblock Cambridge University Press, 1990.

\bibitem[Posner \& Weyl(2018)]{posner2018radical}
Posner, E.A. and Weyl, E.G.
\newblock \emph{Radical Markets}.
\newblock Princeton University Press, 2018.

\bibitem[SourceCred(2020)]{sourcecred2020}
SourceCred.
\newblock A reputation protocol for open communities.
\newblock \url{https://sourcecred.io}, 2020.

\bibitem[Stephan(2012)]{stephan2012economics}
Stephan, P.
\newblock \emph{How Economics Shapes Science}.
\newblock Harvard University Press, 2012.

\bibitem[UNESCO(2021)]{unesco2021science}
UNESCO.
\newblock UNESCO Science Report 2021.
\newblock UNESCO Publishing, 2021.

\bibitem[Weyl et~al.(2022)]{weyl2022decentralized}
Weyl, E.G., Ohlhaver, P., and Buterin, V.
\newblock Decentralized society: Finding web3's soul.
\newblock \emph{SSRN}, 2022.

\bibitem[Wright(2021)]{wright2021daos}
Wright, A.
\newblock The rise of decentralized autonomous organizations.
\newblock \emph{Stanford J.\ Blockchain Law \& Policy}, 4(2):152--237, 2021.

\end{thebibliography}

\end{document}
